\setchapterpreamble[u]{\margintoc}
\chapter{Nature of Light}
\labch{ch:light}

To understand the fundamental physical limits present in optical sensing, it is important to
    understand the physical nature of the processed signal, light. 

Light was perceived in many different ways throughout the history of human kind. 
% TODO: Citation / Reference. I think that this is from the book of science <08-11-23> 
Optics as a field is often considered to have been created by \emph{Ḥasan Ibn al-Haytham}\sidenote{In Europe, he is more
    commonly known by the latinized name \emph{Alhazen}.} in the 10\textsuperscript{th} century. 
% TODO: Citation / Reference <08-11-23> 
There have been many attempts to explain light during the middle ages, but more often than not, they were not in 
    accordance with known experimental properties.
A notable idea was brought forward in particular by \emph{René Descartes}, who described light as a 
    propagating pressure in an unspecified medium\sidenote{This view has interestingly close properties to the way,
    light is described in present day optics.}. % TODO: Reference (unspecified??) <08-11-23> 
The first comprehensive theory of light with regards to the physical reality around us, that was developed in parallel 
    with experimentation, can be dated back to Newton in the 17\textsuperscript{th} century.
In his work, Sir Isaac Newton laid the foundations of the \emph{corpuscular theory of light}\sidecite{newton1704}.
It states that light is composed of particles called \emph{corpuscles} travelling in a straight line with finite
    velocity.
The corpuscular theory replaced in popularity the previous conceptions of Descartes and explained a mayor part of
    conundrums that astronomers had about the optical instrumentation of that time.

% TODO: Finish the history section <08-11-23> 

\section{Wave Theory of Light}%
\labsec{wave-theory-of-light}

% Wave theory of light 
% Maxwells equations 
% Solution to Maxwells equations in the form of a planar wave
% Linearity of the system of equations
% Light in a medium
% Diffraction and the Huygens approximantion
% Fresnel and Frounhoffer diffraction integrals
% Greens function solutions
% Fourier Analysis
%% Spatial-Frequency
% Abbe resolution limit

% Formation of light in a microscope **
% Transfer Functions **
% Transfer function models **
% Principle of SIM **
%% One dimensional case
%% Two dimensional case
% The



% TODO: Note about the paraxial approximation <07-11-23> 
